% year: תשפ"ה
% type: בוחן
% semester: א
% moed: לדוגמה
% official_solution: yes
% source: exams/מבחנים/תשפו/midterm_example 2025.pdf
% part: ב
% number: 4
% points: 10
% categories: functions, multiple-choice

\begin{question}
תהי $f(x)$ פונקציה זוגית. אילו מהטענות הבאות נכונה?
\begin{enumerate}
  \item הפונקציה $g(x)=x^2f(x)-\frac{1}{|x|+1}$ היא פונקציה זוגית.
  \item הפונקציה $g(x)=x^3f(x)-x|x|$ היא פונקציה זוגית.
  \item הפונקציה $g(x)=x^3f(x)+1$ היא פונקציה זוגית.
  \item הפונקציה $g(x)=\frac{f(x)}{x^2+x}$ היא פונקציה אי-זוגית.
\end{enumerate}
\end{question}

\begin{hint}
מכפלה של זוגית בזוגית היא זוגית; מכפלה של אי-זוגית בזוגית היא אי-זוגית. להפרכה מספיקה דוגמה אחת, למשל $f\equiv1$.
\end{hint}

\begin{solution}
\textbf{התשובה הנכונה: (א)}.

\textbf{(א) נכונה:} תחום ההגדרה של $g$ הוא תחום ההגדרה של $f$, שהוא סימטרי (כי $f$ זוגית). לכל $x$ בתחום, מכיוון ש-$f(-x)=f(x)$, $(-x)^2=x^2$ ו-$|-x|=|x|$:
\[
g(-x)=(-x)^2f(-x)-\frac{1}{|-x|+1}=x^2f(x)-\frac{1}{|x|+1}=g(x).
\]

\textbf{למה האחרות שגויות} (הטענות צריכות להתקיים לכל $f$ זוגית, ולכן מספיקה דוגמה נגדית אחת; ניקח $f(x)\equiv1$, שהיא זוגית):
\begin{itemize}
  \item (ב) — $g(x)=x^3-x|x|$: $g(2)=8-4=4$ ואילו $g(-2)=-8+4=-4\neq g(2)$. (באופן כללי $g$ זו היא אי-זוגית, כהפרש של שתי פונקציות אי-זוגיות.)
  \item (ג) — $g(x)=x^3+1$: $g(1)=2$ ואילו $g(-1)=0$.
  \item (ד) — $g(x)=\frac{1}{x^2+x}=\frac{1}{x(x+1)}$ מוגדרת ב-$x=1$ אבל לא ב-$x=-1$, ולכן תחום ההגדרה אינו סימטרי ו-$g$ אינה אי-זוגית. (גם ללא שיקול התחום: $g(2)=\frac16$ ו-$g(-2)=\frac{1}{4-2}=\frac12\neq-\frac16=-g(2)$.)
\end{itemize}
\end{solution}
